\documentclass[12pt,a4paper]{report}

\usepackage{fontspec}
\usepackage{polyglossia}
\setmainlanguage{romanian}
\setmainfont{Latin Modern Roman}

\usepackage[a4paper,margin=2.5cm]{geometry}
\usepackage{setspace}
\usepackage{amsmath}
\usepackage{amssymb}
\usepackage{enumitem}
\usepackage{longtable}
\usepackage{array}
\usepackage{booktabs}
\usepackage{xcolor}
\usepackage{listings}
\usepackage[hidelinks]{hyperref}

\onehalfspacing
\setlength{\parindent}{1.2cm}
\setlength{\parskip}{0.25em}

\definecolor{codebg}{RGB}{248,248,248}
\definecolor{codeframe}{RGB}{210,210,210}
\definecolor{codekeyword}{RGB}{0,78,130}
\definecolor{codestring}{RGB}{120,70,0}
\definecolor{codecomment}{RGB}{80,120,80}

\lstdefinestyle{pythonstyle}{
    language=Python,
    basicstyle=\ttfamily\small,
    keywordstyle=\color{codekeyword}\bfseries,
    stringstyle=\color{codestring},
    commentstyle=\color{codecomment}\itshape,
    backgroundcolor=\color{codebg},
    frame=single,
    rulecolor=\color{codeframe},
    numbers=left,
    numberstyle=\tiny,
    numbersep=8pt,
    breaklines=true,
    showstringspaces=false,
    tabsize=4,
    keepspaces=true
}

\newcommand{\codechunk}[2]{
    \lstinputlisting[
        style=pythonstyle,
        firstline=#1,
        lastline=#2,
        firstnumber=#1
    ]{../src/dsp/vad_energy_zcr.py}
}

\title{Explicatie detaliata a fisierului \texttt{src/dsp/vad\_energy\_zcr.py}}
\author{Documentatie tehnica pentru proiectul de procesare audio}
\date{\today}

\begin{document}

\maketitle
\tableofcontents

\chapter{Rolul fisierului in proiect}

Fisierul \texttt{src/dsp/vad\_energy\_zcr.py} implementeaza detectorul de activitate vocala al proiectului. VAD vine de la \emph{Voice Activity Detection} si reprezinta procesul prin care un semnal audio este impartit in cadre etichetate ca vorbire sau non-vorbire.

Acest modul este important deoarece multe componente ale unui sistem de procesare audio au nevoie sa stie unde exista voce:

\begin{itemize}
    \item in pregatirea datelor, pentru generarea de etichete;
    \item in evaluare, pentru compararea etichetelor prezise cu cele de referinta;
    \item in training, cand scorurile VAD pot fi folosite ca masca sau poarta soft;
    \item in demonstratii real-time, unde sistemul poate afisa sau folosi activitatea vocala cadru cu cadru.
\end{itemize}

Detectorul nu se bazeaza pe un model neuronal. El combina trasaturi clasice de semnal: energie, rata trecerilor prin zero, raport de energie in banda vocala si flatness spectral. Apoi transforma aceste trasaturi intr-un scor adaptiv si il decodeaza cu praguri diferite pentru pornirea si oprirea vorbirii.

\chapter{Ideea algoritmului}

Semnalul audio este impartit in cadre scurte. Pentru fiecare cadru se calculeaza un set de trasaturi:

\[
E_t = \frac{1}{N}\sum_{n=0}^{N-1} x_t[n]^2,
\]

unde \(E_t\) este energia cadrului \(t\), iar \(N\) este lungimea cadrului.

Rata trecerilor prin zero este:

\[
ZCR_t = \frac{1}{N-1}\sum_{n=1}^{N-1}
\mathbf{1}\{sign(x_t[n]) \neq sign(x_t[n-1])\}.
\]

Raportul de banda vocala compara energia din intervalul frecventelor tipice pentru vorbire cu energia spectrala totala:

\[
R_t =
\frac{\sum_{f \in [300,3400]} |X_t(f)|}
{\sum_f |X_t(f)| + \varepsilon}.
\]

Flatness-ul spectral masoara cat de asemanator este spectrul cu zgomotul alb:

\[
F_t =
\frac{
\exp\left(\frac{1}{K}\sum_k \log(|X_t(k)|+\varepsilon)\right)}
{\frac{1}{K}\sum_k (|X_t(k)|+\varepsilon)}.
\]

Un cadru vocal tinde sa aiba energie mai mare, raport bun in banda 300--3400 Hz, ZCR controlat si flatness mai mic decat zgomotul uniform. Fisierul combina aceste observatii intr-un scor:

\[
S_t =
w_E S_E(t) + w_R S_R(t) + w_Z S_Z(t) + w_F S_F(t).
\]

Scorul este apoi transformat in etichete binare prin histerezis: exista un prag mai mare pentru intrarea in vorbire si un prag mai mic pentru iesirea din vorbire. Aceasta alegere evita alternarea rapida intre 0 si 1.

\chapter{Docstring si importuri}

\section{Liniile 1--17}

\codechunk{1}{17}

\begin{description}[leftmargin=2.9cm,style=nextline]
    \item[Liniile 1--10] Contin docstring-ul modulului. El spune ca detectorul foloseste energie, ZCR si indicii spectrale. Mai precizeaza ca API-ul public \texttt{EnergyZCRVAD} este pastrat pentru compatibilitate, chiar daca detectorul este mai robust decat o versiune simpla energie-ZCR.
    \item[Linia 12] Importa \texttt{dataclass}, folosit pentru definirea unei clase de configurare clare si compacte.
    \item[Linia 13] Importa \texttt{numpy}, biblioteca principala pentru calcule numerice pe semnale, cadre, spectre si vectori de etichete.
    \item[Linia 15] Importa \texttt{config}, de unde sunt luati parametrii globali ai proiectului, cum ar fi lungimea cadrului, hop-ul si rata de esantionare.
    \item[Linia 16] Importa \texttt{frame\_signal}, functia comuna care imparte semnalul audio in cadre.
    \item[Linia 17] Importa functia de evaluare VAD din \texttt{metrics}, dar o redenumeste in \texttt{eval\_vad\_metrics}. Redenumirea evita confuzia cu functia locala \texttt{evaluate\_vad\_on\_pair}.
\end{description}

\chapter{Configurarea detectorului}

\section{Liniile 20--44}

\codechunk{20}{44}

Clasa \texttt{EnergyZCRVADConfig} grupeaza toti parametrii detectorului. Fiind decorata cu \texttt{@dataclass}, clasa primeste automat un constructor care accepta aceste campuri ca argumente.

\begin{description}[leftmargin=2.9cm,style=nextline]
    \item[Linia 20] Marcheaza clasa drept dataclass.
    \item[Linia 21] Defineste clasa de configurare.
    \item[Linia 22] Descrie scopul clasei: configurarea detectorului VAD imbunatatit.
    \item[Liniile 23--24] Preiau lungimea cadrului si hop-ul din configuratia globala. Astfel, VAD-ul ramane aliniat cu restul pipeline-ului audio.
    \item[Liniile 26--30] Pastreaza parametri compatibili cu o abordare traditionala: prag relativ de energie, limita ZCR pentru vorbire, numar minim de cadre de vorbire si numar minim de cadre de liniste.
    \item[Liniile 32--34] Definesc pragurile de histerezis. \texttt{speech\_start\_threshold} este mai mare decat \texttt{speech\_end\_threshold}, ceea ce stabilizeaza tranzitiile.
    \item[Liniile 35--38] Definesc ponderile trasaturilor in scorul final: energie, raport de banda, ZCR si flatness.
    \item[Linia 39] Stabileste un scor minim de energie pentru ca un cadru sa poata fi considerat vorbire.
    \item[Liniile 40--41] Definesc banda vocala folosita pentru raportul spectral, de la 300 Hz la 3400 Hz.
    \item[Linia 42] Stabileste limita de flatness peste care cadrul seamana prea mult cu zgomotul.
    \item[Linia 43] Stabileste numarul de cadre \emph{hangover}, adica persistenta scurta a starii de vorbire dupa ce scorul scade.
    \item[Linia 44] Stabileste dimensiunea filtrului median folosit pentru netezirea etichetelor.
\end{description}

Aceasta configurare face detectorul usor de ajustat fara modificarea logicii interne. Pentru experimente se poate crea o instanta \texttt{EnergyZCRVADConfig} cu praguri diferite.

\chapter{Functii de netezire si curatare a etichetelor}

\section{Eliminarea insulelor scurte: liniile 47--62}

\codechunk{47}{62}

Functia \texttt{remove\_short\_islands} elimina segmentele prea scurte dintr-un vector binar. De exemplu, daca se cauta valoarea 1 si apare o secventa de un singur cadru de vorbire intre cadre de liniste, functia o poate transforma inapoi in 0.

\begin{description}[leftmargin=2.9cm,style=nextline]
    \item[Linia 47] Defineste functia. \texttt{labels} este vectorul de etichete, \texttt{value} este valoarea segmentului de curatat, iar \texttt{min\_len} este lungimea minima acceptata.
    \item[Linia 49] Copiaza vectorul pentru a nu modifica direct intrarea primita.
    \item[Liniile 50--51] Initializeaza lungimea vectorului si indexul de parcurgere.
    \item[Liniile 52--61] Parcurg vectorul si cauta segmente consecutive egale cu \texttt{value}.
    \item[Liniile 54--56] Extind indexul \texttt{j} pana la finalul segmentului curent.
    \item[Liniile 57--58] Daca segmentul este mai scurt decat minimul cerut, segmentul este inversat in \texttt{1 - value}.
    \item[Linia 62] Returneaza etichetele curatate.
\end{description}

\section{Filtrul median majoritar: liniile 65--77}

\codechunk{65}{77}

Functia \texttt{median\_smooth} aplica o netezire de tip majoritate pe etichete binare.

\begin{description}[leftmargin=2.9cm,style=nextline]
    \item[Linia 65] Defineste functia pentru vectorul \texttt{labels} si dimensiunea ferestrei.
    \item[Liniile 67--68] Daca fereastra este inutila sau vectorul este gol, functia returneaza intrarea.
    \item[Linia 70] Calculeaza raza ferestrei.
    \item[Linia 71] Adauga margini prin repetarea valorilor de la capete, pentru ca filtrarea sa poata fi aplicata si la inceputul si sfarsitul vectorului.
    \item[Linia 72] Creeaza vectorul rezultat.
    \item[Linia 73] Stabileste pragul majoritar: intr-o fereastra de 5, cel putin 3 valori de 1 produc iesire 1.
    \item[Liniile 74--76] Parcurg fiecare pozitie si aplica regula de majoritate.
    \item[Linia 77] Returneaza vectorul netezit.
\end{description}

\section{Media mobila pentru scoruri: liniile 80--87}

\codechunk{80}{87}

\texttt{moving\_average\_smooth} netezeste scoruri continue, nu etichete binare.

\begin{description}[leftmargin=2.9cm,style=nextline]
    \item[Liniile 80--83] Definesc functia si trateaza cazurile in care nu trebuie aplicata netezire.
    \item[Linia 85] Creeaza un nucleu uniform de mediere.
    \item[Linia 86] Adauga padding la margini pentru a pastra lungimea semnalului.
    \item[Linia 87] Aplica convolutia si returneaza rezultatul ca \texttt{float32}.
\end{description}

\chapter{Calculul trasaturilor de baza}

\section{Energie si ZCR: liniile 90--98}

\codechunk{90}{98}

\begin{description}[leftmargin=2.9cm,style=nextline]
    \item[Linia 90] Defineste functia care intoarce doua siruri: energia si ZCR-ul fiecarui cadru.
    \item[Linia 92] Imparte semnalul in cadre cu \texttt{frame\_signal}.
    \item[Linia 93] Calculeaza energia medie pe fiecare cadru.
    \item[Linia 94] Calculeaza semnul fiecarui esantion.
    \item[Linia 95] Trateaza valorile zero ca valori negative. Astfel se evita ambiguitatea trecerilor prin zero cand un esantion este exact 0.
    \item[Linia 96] Detecteaza schimbarile de semn intre esantioane consecutive.
    \item[Linia 97] Calculeaza rata medie a trecerilor prin zero pentru fiecare cadru.
    \item[Linia 98] Returneaza energia si ZCR-ul.
\end{description}

Energia este utila pentru separarea cadrului vocal de liniste, iar ZCR-ul ajuta la distingerea vorbirii de anumite zgomote sau sunete foarte instabile.

\section{Normalizare robusta: liniile 101--114}

\codechunk{101}{114}

\begin{description}[leftmargin=2.9cm,style=nextline]
    \item[Linia 101] Defineste \texttt{\_safe\_percentile\_range}. Functia calculeaza un interval percentile si il face sigur numeric.
    \item[Liniile 102--103] Calculeaza percentila joasa si percentila inalta.
    \item[Liniile 104--105] Daca intervalul este prea mic, il mareste artificial cu \texttt{eps}. Acest lucru previne impartirea la zero.
    \item[Linia 106] Returneaza limitele.
    \item[Liniile 109--110] Definesc normalizarea unei trasaturi in intervalul \([0,1]\), cu taiere la margini.
    \item[Liniile 113--114] Definesc normalizarea inversata. Aceasta este folosita pentru trasaturi unde valori mici indica mai probabil vorbire, cum ar fi ZCR-ul peste o anumita limita sau flatness-ul.
\end{description}

Normalizarea pe percentile face detectorul adaptiv la nivelul semnalului curent. In loc sa foloseasca praguri absolute fixe pentru fiecare inregistrare, modulul masoara distributia trasaturilor din semnalul primit.

\chapter{Clasa principala \texttt{EnergyZCRVAD}}

\section{Initializare: liniile 117--121}

\codechunk{117}{121}

\begin{description}[leftmargin=2.9cm,style=nextline]
    \item[Linia 117] Defineste clasa principala a detectorului.
    \item[Linia 118] Docstring-ul rezuma metoda: VAD adaptiv cu energie, indicii spectrale si histerezis temporal.
    \item[Linia 120] Constructorul accepta optional o configurare.
    \item[Linia 121] Daca nu se primeste o configurare, se creeaza configurarea implicita.
\end{description}

\section{Decodarea scorurilor in etichete: liniile 123--154}

\codechunk{123}{154}

Metoda \texttt{\_decode\_labels} transforma scorurile continue in etichete binare.

\begin{description}[leftmargin=2.9cm,style=nextline]
    \item[Liniile 123--125] Definesc metoda si initializeaza vectorul de etichete cu zerouri.
    \item[Liniile 126--127] Initializeaza starea curenta si contorul de hangover. Starea 0 inseamna liniste, starea 1 inseamna vorbire.
    \item[Linia 129] Parcurge fiecare scor de cadru.
    \item[Liniile 130--136] Daca detectorul este in starea de liniste, trece in starea de vorbire doar daca scorul depaseste pragul de start si energia este suficienta.
    \item[Liniile 137--147] Daca detectorul este deja in starea de vorbire, acesta iese din vorbire doar cand scorul sau energia scad suficient. Mecanismul de hangover intarzie usor iesirea.
    \item[Linia 149] Scrie starea curenta in vectorul de etichete.
    \item[Linia 151] Aplica netezirea mediana.
    \item[Linia 152] Elimina insulele prea scurte de vorbire.
    \item[Linia 153] Elimina insulele prea scurte de liniste.
    \item[Linia 154] Returneaza etichetele ca intregi pe 32 de biti.
\end{description}

Histerezisul este foarte important in semnale reale. Fara el, cadrele aflate aproape de prag ar produce alternante rapide intre vorbire si liniste.

\chapter{Extragerea trasaturilor}

\section{Liniile 156--194}

\codechunk{156}{194}

Metoda \texttt{compute\_features} calculeaza toate trasaturile folosite ulterior pentru scorul VAD.

\begin{description}[leftmargin=2.9cm,style=nextline]
    \item[Liniile 156--160] Definesc metoda, convertesc semnalul stereo la mono si il transforma in \texttt{float32}.
    \item[Linia 162] Imparte semnalul in cadre.
    \item[Liniile 163--171] Trateaza cazul semnalelor prea scurte sau fara cadre valide, intorcand vectori goi pentru toate trasaturile.
    \item[Linia 173] Creeaza o fereastra Hann.
    \item[Linia 174] Aplica fereastra pe fiecare cadru. Fereastra reduce discontinuitatile de la marginile cadrului in analiza spectrala.
    \item[Linia 176] Calculeaza energia si ZCR-ul prin functia definita anterior.
    \item[Linia 177] Calculeaza log-energia. Logaritmul comprima dinamica energiei si face pragurile mai stabile.
    \item[Linia 179] Calculeaza spectrul de amplitudine prin FFT reala.
    \item[Linia 180] Calculeaza frecventele corespunzatoare binurilor FFT, folosind rata de esantionare din \texttt{config}.
    \item[Linia 181] Creeaza masca pentru banda vocala 300--3400 Hz.
    \item[Linia 183] Calculeaza suma spectrala totala a fiecarui cadru.
    \item[Linia 184] Calculeaza suma spectrala doar in banda vocala.
    \item[Linia 185] Calculeaza raportul de banda vocala.
    \item[Linia 186] Calculeaza flatness-ul spectral, raportul dintre media geometrica si media aritmetica a spectrului.
    \item[Liniile 188--194] Returneaza toate trasaturile intr-un dictionar, convertite la \texttt{float32}.
\end{description}

Aceasta metoda este miezul extractiei de informatii. In loc sa decida doar pe baza energiei, detectorul masoara si structura spectrala a cadrului.

\chapter{Calculul scorului VAD}

\section{Liniile 196--240}

\codechunk{196}{240}

Metoda \texttt{compute\_scores} transforma trasaturile in scoruri de probabilitate relativa pentru vorbire.

\begin{description}[leftmargin=2.9cm,style=nextline]
    \item[Liniile 196--200] Calculeaza trasaturile si trateaza cazul vectorilor goi.
    \item[Liniile 202--205] Calculeaza intervale percentile pentru log-energie, raport de banda, ZCR si flatness.
    \item[Liniile 207--208] Normalizeaza energia si raportul de banda astfel incat valori mai mari sa dea scor mai mare.
    \item[Liniile 210--211] Stabileste limita superioara pentru ZCR si calculeaza scorul inversat. ZCR prea mare este penalizat.
    \item[Liniile 213--214] Stabileste limita pentru flatness si calculeaza scorul inversat. Flatness mare inseamna spectru mai apropiat de zgomot.
    \item[Liniile 216--222] Creeaza vectorul de ponderi si il normalizeaza ca suma sa fie 1.
    \item[Liniile 224--229] Combina trasaturile normalizate intr-un scor final.
    \item[Liniile 231--232] Aplica o penalizare daca scorul de energie este sub pragul minim relativ.
    \item[Liniile 234--239] Adauga scorurile intermediare in dictionarul de trasaturi. Acest lucru ajuta la inspectie, debug si vizualizare.
    \item[Linia 240] Returneaza scorul final si dictionarul extins de trasaturi.
\end{description}

Scorul final nu trebuie citit ca probabilitate calibrata statistic. El este mai degraba o masura adaptiva de incredere ca un cadru contine vorbire.

\chapter{Predictii hard si soft}

\section{Predictie binara: liniile 242--256}

\codechunk{242}{256}

Metoda \texttt{predict} produce etichete binare.

\begin{description}[leftmargin=2.9cm,style=nextline]
    \item[Liniile 242--251] Definesc metoda si documenteaza intrarea si iesirea.
    \item[Linia 252] Calculeaza scorurile si trasaturile.
    \item[Liniile 253--254] Daca nu exista scoruri, intoarce un vector gol de etichete.
    \item[Linia 256] Decodeaza scorurile in etichete binare folosind energia ca protectie suplimentara.
\end{description}

\section{Predictie soft: liniile 258--274}

\codechunk{258}{274}

Metoda \texttt{predict\_soft} intoarce un scor continuu in intervalul \([0,1]\).

\begin{description}[leftmargin=2.9cm,style=nextline]
    \item[Liniile 258--265] Documenteaza scopul metodei: scor de incredere per cadru, util in special pentru training sau gating soft.
    \item[Linia 266] Calculeaza scorurile.
    \item[Liniile 267--268] Trateaza cazul fara cadre.
    \item[Linia 270] Calculeaza etichetele binare, apoi le converteste la \texttt{float32}.
    \item[Linia 271] Limiteaza scorul brut in intervalul \([0,1]\).
    \item[Linia 272] Pentru cadrele confirmate ca vorbire, ridica scorul cel putin la pragul de final de vorbire. Astfel, cadrele acceptate de histerezis nu raman cu scor prea mic.
    \item[Linia 273] Netezeste scorul soft prin medie mobila.
    \item[Linia 274] Returneaza scorul final taiat la intervalul valid.
\end{description}

Predictia soft este utila cand nu se doreste o decizie dura 0/1. De exemplu, un model de enhancement poate folosi scorul ca factor de ponderare.

\chapter{Alinierea si evaluarea}

\section{Alinierea VAD: liniile 276--288}

\codechunk{276}{288}

Metoda statica \texttt{align\_vad\_to\_frames} adapteaza lungimea unui vector VAD la un numar dorit de cadre.

\begin{description}[leftmargin=2.9cm,style=nextline]
    \item[Linia 276] Marcheaza metoda ca statica. Ea nu foloseste starea obiectului.
    \item[Linia 277] Defineste metoda cu vectorul VAD si numarul de cadre tinta.
    \item[Liniile 279--280] Daca lungimea este deja corecta, returneaza vectorul.
    \item[Liniile 281--282] Daca vectorul este prea lung, il taie.
    \item[Liniile 283--284] Daca vectorul este gol, intoarce zerouri.
    \item[Liniile 285--287] Daca vectorul este prea scurt, copiaza valorile existente si completeaza finalul cu ultima eticheta disponibila.
    \item[Linia 288] Returneaza vectorul aliniat.
\end{description}

\section{Evaluarea pe o pereche: liniile 291--305}

\codechunk{291}{305}

\begin{description}[leftmargin=2.9cm,style=nextline]
    \item[Liniile 291--295] Definesc functia de evaluare pe un semnal si etichetele lui reale.
    \item[Liniile 296--298] Documenteaza functia.
    \item[Liniile 299--300] Daca nu se primeste un model VAD, creeaza unul implicit.
    \item[Linia 301] Prezice etichetele VAD pentru semnalul audio.
    \item[Linia 302] Alege lungimea comuna minima intre predictii si etichetele reale.
    \item[Liniile 303--304] Taie ambii vectori la aceeasi lungime.
    \item[Linia 305] Apeleaza functia de metrici VAD si returneaza rezultatul.
\end{description}

Aceasta functie conecteaza detectorul la partea de evaluare a proiectului. Ea ascunde detaliile de aliniere minima si returneaza direct dictionarul de metrici.

\chapter{Fluxul complet al modulului}

Fluxul de lucru al fisierului este urmatorul:

\begin{enumerate}
    \item semnalul audio este primit ca tablou NumPy;
    \item daca are mai multe canale, este convertit la mono;
    \item semnalul este impartit in cadre;
    \item pentru fiecare cadru se calculeaza energia, log-energia, ZCR, raportul de banda vocala si flatness-ul spectral;
    \item trasaturile sunt normalizate adaptiv folosind percentile;
    \item scorurile individuale sunt combinate ponderat;
    \item energia prea mica penalizeaza scorul final;
    \item scorul este decodat cu histerezis si hangover;
    \item etichetele sunt netezite si segmentele prea scurte sunt eliminate;
    \item rezultatul final este un vector de etichete 1 pentru vorbire si 0 pentru liniste.
\end{enumerate}

\chapter{De ce este fisierul important}

\texttt{vad\_energy\_zcr.py} este important pentru ca ofera proiectului o componenta VAD explicabila, rapida si independenta de antrenare. Spre deosebire de o retea neuronala, acest detector poate fi inteles direct din formule si praguri.

In acelasi timp, nu este doar un VAD simplu bazat pe energie. El introduce:

\begin{itemize}
    \item normalizare adaptiva pe fiecare semnal;
    \item indicii spectrale specifice vorbirii;
    \item penalizare pentru zgomote spectrale plate;
    \item histerezis pentru tranzitii stabile;
    \item smoothing si eliminarea segmentelor imposibil de scurte.
\end{itemize}

Aceste mecanisme il fac potrivit pentru un proiect de procesare audio in timp real, unde deciziile trebuie sa fie rapide, dar suficient de stabile pentru a nu crea fluctuatii vizuale sau functionale.

\chapter{Concluzie}

Fisierul \texttt{src/dsp/vad\_energy\_zcr.py} implementeaza un detector VAD adaptiv si robust, bazat pe trasaturi clasice de semnal. El combina energia, ZCR-ul, informatia din banda vocala si flatness-ul spectral intr-un scor unic, apoi foloseste histerezis si netezire pentru a produce etichete stabile.

Prin aceasta structura, modulul ocupa un loc central intre procesarea audio clasica si componentele de training/evaluare. Este un fisier esential pentru intelegerea modului in care proiectul detecteaza vorbirea si foloseste aceasta informatie in pipeline-ul audio.

\end{document}
